\documentclass[11pt]{article}
\usepackage[margin=1in]{geometry}
\usepackage{amsmath,amssymb,amsthm}
\usepackage{graphicx}
\usepackage[hidelinks]{hyperref}
\usepackage{enumitem}

\newtheorem{theorem}{Theorem}
\newtheorem{lemma}[theorem]{Lemma}
\newtheorem{corollary}[theorem]{Corollary}
\theoremstyle{definition}
\newtheorem{definition}[theorem]{Definition}

\newcommand{\cU}{\mathcal U}
\newcommand{\cP}{\mathcal P}

\title{An Arbitrary-Family Diagonal Principle for Nets}
\author{Candidate full solution for arXiv:2304.04189}
\date{July 22, 2026}

\begin{document}
\maketitle

\begin{center}
\fbox{\parbox{0.9\textwidth}{\textbf{Status.} Candidate full solution,
likely valid pending expert review.  The proof is in ZFC and explicitly uses
choice; it is not claimed to be constructive.}}
\end{center}

\section{Source question}

Azouzi proves a diagonal principle for a countable sequence of properties of
nets \cite{Azouzi}.  At the start of Section 4, PDF page 7, he asks whether
there is ``a more general diagonal process where the set of properties is not
assumed countable,'' observes that a constructive proof should not be
expected, and suggests Zorn's lemma.  The complete source passage is shown
below.

\begin{center}
\includegraphics[width=\textwidth]{figures/open_problem_crop.png}
\end{center}

The natural precise formulation is obtained from Theorem 7 of the source by
replacing $(\cP_k)_{k\in\mathbb N}$ with a family indexed by an arbitrary set.

\section{Definitions}

We use the subnet convention of the source: $y\colon B\to E$ is a subnet of
$x\colon A\to E$ if $y=x\circ\phi$ for a \emph{final} map
$\phi\colon B\to A$, meaning that for every $\alpha_0\in A$ there is
$\beta_0\in B$ such that $\phi(\beta)\geq\alpha_0$ whenever
$\beta\geq\beta_0$.  It is an \emph{eventual subnet} if some tail is a
subnet.

For a net $x$ in $E$, its eventuality filter is
\[
 \mathfrak F_x=\{S\subseteq E:x\text{ is eventually in }S\}.
\]
The net is \emph{universal} if, for every $S\subseteq E$, it is eventually in
$S$ or eventually in $E\setminus S$.  Equivalently, $\mathfrak F_x$ is an
ultrafilter.

\section{Proof intuition}

Countable diagonalization repeatedly refines a net and then follows one
diagonal through the resulting chain.  For an arbitrary family there is no
last stage and no elementary transfinite limit subnet.  The replacement is to
make the tail behavior maximal first.

Pass to a universal subnet $u$.  Every further subnet of $u$ has exactly the
same eventuality ultrafilter.  For each desired property choose a satisfying
subnet $v^i$ of $u$.  Because all their tail ranges belong to one ultrafilter,
every finite collection of tail requirements has a common value.  Indexing
over all finite collections of requirements produces one net which is a
subnet of $u$ and, on a tail, a subnet of every $v^i$.  Heredity under
eventual subnets then transfers all properties at once.

\section{Two lemmas}

\begin{lemma}[Universal subnet]
Every net has a universal subnet, assuming the ultrafilter lemma.
\end{lemma}

\begin{proof}
Let $x\colon A\to E$ and put
$T_\alpha=\{x_\gamma:\gamma\geq\alpha\}$.  The sets $T_\alpha$ form a
filter base.  Extend the generated filter to an ultrafilter $\cU$ on $E$.

Let
\[
 D=A\times\cU,
 \qquad
 (\alpha,V)\preceq(\alpha',V')
 \quad\Longleftrightarrow\quad
 \alpha\leq\alpha'\ \text{ and }\ V'\subseteq V.
\]
This is directed: use an upper bound in $A$ and the intersection of the two
ultrafilter sets.  Since $T_\alpha\cap V\neq\varnothing$, choose
$\psi(\alpha,V)\geq\alpha$ with $x_{\psi(\alpha,V)}\in V$, and define
$u_{(\alpha,V)}=x_{\psi(\alpha,V)}$.  The map $\psi\colon D\to A$ is final,
so $u$ is a subnet of $x$.

For $S\subseteq E$, either $S\in\cU$ or $E\setminus S\in\cU$.  In the first
case, after any index of the form $(\alpha,S)$, every later value of $u$ lies
in $S$; the second case is identical with the complement.  Hence $u$ is
universal.
\end{proof}

\begin{lemma}[Arbitrary common eventual refinement]
Let $(v^j\colon B_j\to E)_{j\in J}$ be a nonempty family of nets with the same
eventuality ultrafilter $\cU$.  Fix $j_*\in J$.  There is a net $w$ which is a
subnet of $v^{j_*}$ and an eventual subnet of $v^j$ for every $j\in J$.
\end{lemma}

\begin{proof}
For $b\in B_j$, write
\[
 R_j(b)=\{v^j_c:c\geq b\}.
\]
Because $v^j$ is eventually in its own tail range, $R_j(b)\in\cU$.

Let $D$ consist of pairs
\[
 d=(F,(b_j)_{j\in F}),
\]
where $F\subseteq J$ is finite, $j_*\in F$, and $b_j\in B_j$.  Order $D$ by
\[
 (F,(b_j))\preceq(G,(c_j))
 \quad\Longleftrightarrow\quad
 F\subseteq G\ \text{ and }\ b_j\leq c_j\quad(j\in F).
\]
The set $D$ is directed: for two constraints take the union of their finite
sets and, on every overlapping coordinate, an upper bound in $B_j$.

For every $d\in D$, the finite intersection
\[
 Q_d=\bigcap_{j\in F}R_j(b_j)
\]
belongs to $\cU$ and is nonempty.  Choose $w_d\in Q_d$.  Fix $j\in J$ and
let $D_j=\{d\in D:j\in F_d\}$.  This is a tail of $D$.  For $d\in D_j$,
choose $\theta_j(d)\in B_j$ such that
\[
 \theta_j(d)\geq b_j(d)
 \quad\text{and}\quad
 v^j_{\theta_j(d)}=w_d.
\]
The map $\theta_j\colon D_j\to B_j$ is final.  Indeed, for $b_0\in B_j$,
choose a constraint $d_0$ containing $j$ with $b_j(d_0)=b_0$.  Whenever
$d\succeq d_0$, one has
$\theta_j(d)\geq b_j(d)\geq b_0$.  Thus $w|_{D_j}$ is a subnet of $v^j$.

Since $j_*$ belongs to every $F_d$, $D_{j_*}=D$; hence $w$ itself is a subnet
of $v^{j_*}$.  For every other $j$, the same conclusion holds on the tail
$D_j$, so $w$ is an eventual subnet of $v^j$.
\end{proof}

\section{Arbitrary-family diagonal principle}

\begin{theorem}
Let $E$ be a nonempty set, let $x\colon A\to E$ be a net, and let
$(\cP_i)_{i\in I}$ be a family of properties of nets indexed by an arbitrary
set $I$.  Assume:
\begin{enumerate}[label=\textup{(\roman*)}]
\item each $\cP_i$ is preserved under passing to eventual subnets;
\item for every $i\in I$ and every subnet $y$ of $x$, there is a further
subnet of $y$ satisfying $\cP_i$.
\end{enumerate}
Then $x$ has a subnet satisfying $\cP_i$ for every $i\in I$.
\end{theorem}

\begin{proof}
The conclusion is immediate if $I=\varnothing$, so assume $I\neq\varnothing$.
By the universal-subnet lemma, choose a universal subnet $u$ of $x$, and let
$\cU=\mathfrak F_u$.

For each $i\in I$, hypothesis (ii), applied to $u$, gives a subnet $v^i$ of
$u$ satisfying $\cP_i$.  Every subnet of a universal net is universal: it is
eventually in whichever of $S$ and $E\setminus S$ eventually contains $u$.
Moreover, $\mathfrak F_{v^i}\supseteq\mathfrak F_u=\cU$.  Both sides are
ultrafilters, so
\[
 \mathfrak F_{v^i}=\cU\qquad(i\in I).
\]

Apply the common-refinement lemma to the family consisting of $u$ and all
$v^i$, taking $u$ as the distinguished member.  We obtain a net $z$ which is
a subnet of $u$ and an eventual subnet of every $v^i$.  Since $v^i$ satisfies
$\cP_i$ and $\cP_i$ is preserved under eventual subnets, $z$ satisfies
$\cP_i$ for every $i\in I$.  Finally, the composition of final maps is final,
so the subnet $z$ of $u$ is a subnet of $x$.
\end{proof}

\section{Consequences}

\begin{corollary}[Tychonoff]
Assuming Choice, an arbitrary product of compact spaces is compact.
\end{corollary}

\begin{proof}
Let $x$ be a net in $\prod_{i\in I}X_i$.  For each $i$, let $\cP_i$ mean
that the $i$th coordinate net converges.  Compactness of $X_i$ says that every
subnet has a further subnet satisfying $\cP_i$, and coordinate convergence is
preserved under subnets.  The theorem produces one subnet converging in every
coordinate, hence in the product topology.
\end{proof}

\begin{corollary}[Banach--Alaoglu]
For every Banach space $X$, the dual unit ball is weak-star compact.
\end{corollary}

\begin{proof}
Given a net in $B_{X^*}$ and $x\in X$, let $\cP_x$ mean that the scalar
evaluation net at $x$ converges.  Every subnet has a further subnet satisfying
$\cP_x$, because its values lie in a compact scalar disk.  The theorem gives
a subnet converging pointwise on all of $X$.  Its pointwise limit is linear
and bounded by one, hence belongs to $B_{X^*}$, and the subnet converges to it
in the weak-star topology.
\end{proof}

\section{Verification and limitations}

The proof matches the final-map definition of subnet used in the source.  In
the common-refinement lemma, the distinguished component is refined on the
entire index set; each other component is refined on a tail, exactly matching
the heredity hypothesis.  No countability assumption enters.

The result is nonconstructive.  It uses the ultrafilter lemma and additional
set-indexed selections, all available in ZFC.  No assertion is made about the
corresponding statement in ZF alone.

The run's lightweight indexes and a bounded web search on July 22, 2026 were
checked using arXiv:2304.04189, the exact title and author, phrases from the
source question, and close variants involving universal nets and common
ultrafilters.  No later explicit resolution was found.  Since universal nets
are classical, the exact theorem may nevertheless be folklore; novelty
confidence is moderate pending specialist review.

\begin{thebibliography}{9}
\bibitem{Azouzi}
Y.~Azouzi,
\emph{A diagonal principle for nets},
arXiv:2304.04189 (2023).

\bibitem{Kelley}
J.~L.~Kelley,
\emph{General Topology},
Van Nostrand, 1955.
\end{thebibliography}

\end{document}
